% Written for Swift Shifter's benchmark (#150); part of the repository,
% under its licence (GPL-3.0-or-later).
\documentclass[11pt]{ctexart}
\usepackage{amsmath,amssymb,amsthm}

\newtheorem{theorem}{定理}[section]
\newtheorem{lemma}[theorem]{引理}
\theoremstyle{definition}
\newtheorem{definition}[theorem]{定義}

\newcommand{\R}{\mathbb{R}}
\newcommand{\norm}[1]{\left\lVert #1 \right\rVert}
\DeclareMathOperator{\tr}{tr}

\title{矩陣範數與 Gradient Descent 的收斂}
\author{Swift Shifter 測試文件}
\date{2026 年 10 月}

\begin{document}
\maketitle

\begin{abstract}
本文以中英混排的方式，整理 matrix norm 的幾個基本性質，並說明 gradient
descent 在 strongly convex 函數上的線性收斂。文中的公式包含行內公式、編號公式、
多行對齊與定理環境，用來測試轉換器是否保留數學內容。
\end{abstract}

\section{矩陣範數}

對矩陣 $A \in \R^{m \times n}$，Frobenius norm 定義為
\begin{equation}
  \norm{A}_F = \sqrt{\sum_{i=1}^{m} \sum_{j=1}^{n} a_{ij}^2} = \sqrt{\tr(A^\top A)}.
  \label{eq:frobenius}
\end{equation}
由式~\eqref{eq:frobenius} 可知 $\norm{A}_F = \norm{A^\top}_F$。另一個常用的範數是
spectral norm $\norm{A}_2 = \sigma_{\max}(A)$，也就是最大的 singular value。

\begin{definition}
若對所有 $x, y \in \R^n$ 都有 $f(y) \ge f(x) + \nabla f(x)^\top (y - x) + \frac{\mu}{2}\norm{y - x}^2$，
則稱 $f$ 為 $\mu$-strongly convex。
\end{definition}

\begin{lemma}\label{lem:norms}
對任意 $A \in \R^{m \times n}$，若 $r = \operatorname{rank}(A)$，則
\[
  \norm{A}_2 \le \norm{A}_F \le \sqrt{r}\, \norm{A}_2 .
\]
\end{lemma}

\begin{proof}
設 $\sigma_1 \ge \dots \ge \sigma_r > 0$ 為 $A$ 的 singular values，則
$\norm{A}_F^2 = \sum_{k=1}^{r} \sigma_k^2$，兩邊的不等式都由此直接得到。
\end{proof}

\section{Gradient Descent 的收斂}

考慮更新規則
\begin{align}
  x_{k+1} &= x_k - \eta \nabla f(x_k), \label{eq:gd}\\
  e_k     &= f(x_k) - f(x^\star). \nonumber
\end{align}

\begin{theorem}\label{thm:linear}
若 $f$ 為 $\mu$-strongly convex 且 $\nabla f$ 為 $L$-Lipschitz，取步長 $\eta = 1/L$，
則依式~\eqref{eq:gd} 更新時
\begin{equation}
  e_{k} \le \left(1 - \frac{\mu}{L}\right)^{k} e_0 .
\end{equation}
\end{theorem}

定理~\ref{thm:linear} 說明誤差以幾何級數下降，比值 $\kappa = L/\mu$ 稱為
condition number。$\kappa$ 越大，收斂越慢；這也是為什麼實務上常用
preconditioning 或 momentum。引理~\ref{lem:norms} 在估計 $L$ 時很有用：
對二次函數 $f(x) = \frac12 x^\top A x$，$L = \norm{A}_2 \le \norm{A}_F$。

\end{document}
